\chapter{2022 Nobel Prize in Economics}
\textbf{Laureates: }Ben Bernanke (USA), Douglas Diamond (USA), Philip Dybvig (USA)

\section*{Reason for Award}
Research on banks and financial crises

\section{What They Did}
Ben Bernanke, Douglas Diamond, and Philip Dybvig advanced understanding of banks and financial crises. Bernanke studied the causes and consequences of the Great Depression, showing that bank failures exacerbated the crisis by destroying credit intermediation. He argued that deposit insurance and lender of last resort functions prevent bank runs. Diamond and Dybvig developed the model of bank runs, showing how banks provide liquidity transformation by accepting short-term deposits and making long-term loans. Their model demonstrated that bank runs can occur as self-fulfilling equilibria when depositors fear others will withdraw first. They showed how deposit insurance and capital requirements can prevent runs, while also creating moral hazard problems.

\section{Impact on Economic Thinking}
Bernanke, Diamond, and Dybvig transformed understanding of financial intermediation and crisis dynamics. Bernanke's historical analysis influenced monetary policy and financial regulation, emphasizing the importance of preventing bank failures. Diamond-Dybvig model became the standard framework for analyzing bank runs, liquidity provision, and financial stability. Their work explained why banks are both essential and fragile, providing foundations for banking regulation. The model informed discussions about deposit insurance, bailouts, and resolution mechanisms. Their research highlighted the trade-off between financial stability and moral hazard, influencing regulatory design. The framework has been extended to analyze shadow banking, money market funds, and cryptocurrency risks.

\section{Modern Perspectives Today}
Banking theory based on Diamond-Dybvig and Bernanke's work remains essential for understanding financial stability. The 2008 financial crisis validated their insights about liquidity transformation, contagion, and systemic risk. Modern banking regulation, including capital requirements, stress tests, and resolution frameworks, draws on their theoretical foundations. Research on shadow banking, money market funds, and central bank digital currencies extends their framework to new contexts. The COVID-19 pandemic highlighted the importance of bank liquidity and depositor confidence. Climate risk and green finance create new questions about bank stability that build on their work. As financial innovation continues, Diamond, Dybvig, and Bernanke's insights remain relevant for understanding and preventing future crises.
